\chapter{Literature Review}
\label{ch:lit}

This chapter reviews the work on which MandiSense AI rests and states, from that work, the problem the thesis addresses. The review has two parts. The first is a systematic survey, carried out in the initial stage of the project, of \textbf{25 recent studies} on agricultural price prediction, volatility and anomaly analysis, hybrid and ensemble forecasting, and smart-farming decision support (Section~\ref{sec:survey}); each study is compared on its techniques, results, research gaps and takeaways. The second widens the view to the methods the system uses (decision-focused learning, conformal prediction, missing data, forecast evaluation and language-model grounding) and to the economics of Indian agricultural markets. Every reference in the Bibliography was looked up by title and author in a public bibliographic index (Crossref or arXiv) before it was listed.

Section~\ref{sec:lit-markets} describes the setting and the evidence that information changes farmer outcomes. Sections~\ref{sec:related-forecast} to~\ref{sec:evaluation} cover forecasting, decomposition hybrids and evaluation. Section~\ref{sec:survey} presents the survey with comparison tables. Sections~\ref{sec:ml-decision} to~\ref{sec:llm} cover decisions, uncertainty, missing data, transmission and language models. Section~\ref{sec:limitations} gathers the gaps and Section~\ref{sec:problem} states the problem. Figure~\ref{fig:taxonomy} places the strands of work and the point at which this project joins them.

\begin{figure}[!b]
\centering
\begin{tikzpicture}[node distance=6mm, every node/.style={font=\small}, box/.style={draw, rounded corners=3pt, align=center, minimum height=9mm, text width=3.6cm, fill=blue!6}, hi/.style={box, fill=orange!20, very thick}]
\node[box] (a) {Price forecasting\\ ARIMA, LSTM, hybrids};
\node[box, right=of a] (b) {Decomposition hybrids\\ STL, VMD, ICEEMDAN};
\node[box, right=of b] (c) {Foundation models\\ Chronos, zero-shot};
\node[box, below=of a] (d) {Decision-focused learning\\ SPO+, benchmarks};
\node[box, below=of b] (e) {Conformal prediction\\ adaptive, delayed feedback};
\node[box, below=of c] (f) {Missing data\\ not-at-random reporting};
\node[hi, below=10mm of e] (g) {\textbf{MandiSense AI}\\ outcome-supervised market choice\\ with verified floors};
\draw[->] (a) -- (d); \draw[->] (b) -- (e); \draw[->] (c) -- (e); \draw[->] (f) -- (g); \draw[->] (d) -- (g); \draw[->] (e) -- (g);
\end{tikzpicture}
\caption{Strands of related work and where this project joins them.}
\label{fig:taxonomy}
\end{figure}

\section{Agricultural Markets and the Value of Information}
\label{sec:lit-markets}

Fruit and vegetable prices in India are volatile, and the volatility is not only a matter of weather: a study of onion prices traces large swings to supply shocks interacting with storage, trade policy and market structure~\cite{birthal2019}. The government publishes daily mandi quotations through AGMARKNET~\cite{agmarknet}, but a quotation is the price at one market on one day, not advice on where to sell.

The effect of price information on producers has been studied in several countries. In the South Indian fisheries sector, the arrival of mobile phones reduced price dispersion and raised both producer profit and consumer surplus~\cite{jensen2007}. In Niger, phones reduced search costs and price dispersion across grain markets~\cite{aker2010}. For Indian agriculture, a randomised study of text-message price information found that it changed some marketing decisions but had a limited average effect~\cite{fafchamps2012}.

\paragraph{What this means for the project.} The studies agree that a price quotation helps when it changes where or when a farmer sells, and that the effect is small when the information is generic or late. MandiSense AI therefore delivers a specific choice, in the farmer's language, with the date of every price shown (Chapter~\ref{ch:impl}), and not a price board. We do not claim any welfare effect: that would need a field trial, which is listed as future work (Chapter~\ref{ch:concl}).

\section{Related Work on Agricultural Price Forecasting}
\label{sec:related-forecast}

\subsection*{Statistical and neural forecasters}
Forecasting texts treat the random walk and seasonal methods as the benchmarks any new model must beat~\cite{hyndman2021}. Early work on Indian crops used time-delay neural networks for monthly wholesale prices of oilseeds~\cite{jha2014}. Later studies combined linear and recurrent models, for example an ARIMA--LSTM hybrid with random-forest feature selection for volatile series~\cite{ray2023}; applied deep models to ragi prices in Karnataka~\cite{meena2024}; guided networks with hidden Markov regimes for highly volatile commodities~\cite{avinash2024}; and, for multi-commodity settings, coupled attention models with anomaly detection~\cite{sediyono2025}. A modular machine-learning framework predicts crop yield, supply and demand before the price~\cite{mohanty2023}, and deep models with exogenous drivers such as weather and supply factors improved forecasts of tomato, onion and potato prices~\cite{exogenous2024}. A recent review catalogues this architectural variety and flags open challenges~\cite{wu2026}. Work on farmers' choices shows that crop insurance affects diversification, which links price information to planting decisions~\cite{crop_ins}.

\subsection*{Volatility}
Volatility is the quantity a farmer feels as risk. A hybrid LSTM--GARCH model, applied to 23 commodities across 165 Indian markets, captured both non-linear dynamics and volatility clustering and gave lower error than standalone econometric and deep models~\cite{manogna2025}. This is the model family behind the volatility view of the trader desk (Chapter~\ref{ch:impl}).

\subsection*{Gradient boosting and foundation models}
Gradient-boosted trees~\cite{chen2016,ke2017} remain strong on tabular data, where they often outperform deep networks~\cite{grinsztajn2022}. Pretrained time-series foundation models forecast new series without training~\cite{ansari2024}, and early evidence for commodity prices is cautiously positive~\cite{wang2026}. Both belong in a fair benchmark, so both are included in Chapter~\ref{ch:results}.

\section{Decomposition Hybrids and the Leakage Problem}
\label{sec:hybrids}

A large strand of work splits the price into components and forecasts each part. The classical decomposition is STL, which separates trend, seasonal and remainder by loess~\cite{cleveland1990}; variational mode decomposition extracts intrinsic modes by constrained optimisation~\cite{dragomiretskiy2014}. Reported gains are large: an STL--LSTM hybrid for seasonal series~\cite{jaiswal2025}, CEEMDAN-based hybrids for non-linear prices~\cite{ceemdan2024}, ICEEMDAN--VMD--GRU~\cite{sathe2025}, ensembles of machine-learning, time-series and stochastic models~\cite{celik2025}, and a decomposition-ensemble with kernel density estimation for point and interval forecasts~\cite{kde2024}.

The concern with this design is that a decomposition computed on the whole series before the split lets future values shape the components used to predict them. The evaluation literature warns about time-ordered validation in general~\cite{bergmeir2012,cerqueira2020}, but we found no agricultural study that measured the size of the leak. Section~\ref{sec:leakage} does.

\section{Evaluation and Statistical Comparison}
\label{sec:evaluation}

Predictive accuracy is compared pairwise with the Diebold--Mariano test~\cite{diebold1995}, with a small-sample correction~\cite{harvey1997} and autocorrelation-robust variance~\cite{newey1987}. When many comparisons are made, familywise error is controlled by Holm's procedure~\cite{holm1979}, and a model confidence set identifies the models that cannot be separated from the best~\cite{hansen2011}. Dependence in the data is respected by resampling blocks rather than days~\cite{kunsch1989}. Hyper-parameters are tuned inside the training window by Bayesian optimisation~\cite{akiba2019}, and the test period is touched once.

\section{Systematic Survey of 25 Studies (2023--2025)}
\label{sec:survey}

\subsection{Survey protocol}
The survey was carried out in the initial stage of the project and recorded in a shared workbook with the same fields for every paper: title, publication, abstract, techniques, results, year, link, research gaps and key takeaways. The work was divided among the three team members. The search was limited to peer-reviewed journal and conference papers from 2023 to 2025 on agricultural price prediction, volatility and anomaly analysis, hybrid and ensemble forecasting, and smart-farming decision support, drawn from Springer Nature, Elsevier ScienceDirect and IEEE Xplore. After removing duplicates (six IEEE papers were recorded by two members), \textbf{25 distinct studies} remained, labelled S1 to S25. Every study was matched to a bibliographic record in Crossref before it was cited. The findings below are the authors' own, as recorded in the workbook from the abstracts and results sections; performance figures were not re-measured by us.

\begin{table}[!t]
\centering
\caption{Distribution of the 25 surveyed studies by publisher and year (as recorded in the workbook)}
\label{tab:survey-dist}
\small
\begin{tabular}{|l|c|c|c|c|}
\hline
\textbf{Source} & \textbf{2023} & \textbf{2024} & \textbf{2025} & \textbf{Total} \\ \hline
Elsevier (ScienceDirect), S1--S10 & 2 & 5 & 3 & 10 \\ \hline
Springer Nature, S11--S19 & 2 & 5 & 2 & 9 \\ \hline
IEEE, S20--S25 & 1 & 3 & 2 & 6 \\ \hline
\textbf{Total} & \textbf{5} & \textbf{13} & \textbf{7} & \textbf{25} \\ \hline
\end{tabular}
\end{table}

\subsection{Overview of the studies}
Table~\ref{tab:survey-overview} lists the studies with their venue and the task they address. Of the 25, 15 forecast prices (S2 also detects anomalies), 5 concern volatility or its drivers, 1 is about anomaly detection, and 4 are crop-recommendation systems or reviews.

\Needspace{12\baselineskip}
{\footnotesize\setlength{\tabcolsep}{4pt}\renewcommand{\arraystretch}{1.15}
\begin{longtable}{|p{0.9cm}|p{7.7cm}|p{4.2cm}|p{2.7cm}|}
\caption{The 25 surveyed studies}\label{tab:survey-overview}\\
\hline \textbf{ID} & \textbf{Study} & \textbf{Venue} & \textbf{Task} \\ \hline \endfirsthead
\hline \textbf{ID} & \textbf{Study} & \textbf{Venue} & \textbf{Task} \\ \hline \endhead
\hline \endfoot
S1 & Hybrid forecasting of commodity prices: machine learning, time series and stochastic simulation~\cite{celik2025} & Borsa \.Istanbul Rev., 2025 & Price, risk \\ \hline
S2 & Multi-commodity price forecasting and anomaly detection with attention-boosted models~\cite{sediyono2025} & J. Agric. Food Res., 2025 & Price, anomaly \\ \hline
S3 & ARIMA--LSTM for volatile prices with random forest~\cite{ray2023} & Appl. Soft Comput., 2023 & Price \\ \hline
S4 & ICEEMDAN--VMD--GRU for economic resilience~\cite{sathe2025} & Results Eng., 2025 & Price \\ \hline
S5 & Volatility dynamics of agricultural futures under uncertainties~\cite{fiszeder2024} & Energy Econ., 2024 & Volatility \\ \hline
S6 & Financial stress and realised volatility of agricultural commodities~\cite{algirdas2024} & Res. Int. Bus. Finance, 2024 & Volatility \\ \hline
S7 & Optimal forecast combination (PSO-CS) for daily futures prices~\cite{zeng2023} & Appl. Soft Comput., 2023 & Price, ensemble \\ \hline
S8 & Transformer anomaly detection with inter-variable attention~\cite{tranad2024} & Inf. Fusion, 2024 & Anomaly \\ \hline
S9 & Hidden-Markov-guided deep models for volatile prices~\cite{avinash2024} & Appl. Soft Comput., 2024 & Price \\ \hline
S10 & Economic policy uncertainty and agricultural prices, China~\cite{epu2024} & Asia Global Econ., 2024 & Volatility \\ \hline
S11 & STL--LSTM for seasonal agricultural price series~\cite{jaiswal2025} & Ann. Data Sci., 2025 & Price \\ \hline
S12 & CEEMDAN hybrid models for agricultural prices~\cite{ceemdan2024} & Sci. Rep., 2024/25 & Price \\ \hline
S13 & Exogenous-variable deep learning for tomato, onion, potato~\cite{exogenous2024} & Sci. Rep., 2024 & Price \\ \hline
S14 & Hybrid LSTM--GARCH volatility forecasting, India~\cite{manogna2025} & J. Big Data, 2025 & Volatility \\ \hline
S15 & Commodity price prediction: a machine-learning framework~\cite{mohanty2023} & Neural Comput. Appl., 2023 & Price \\ \hline
S16 & Demand and price forecasting using deep learning~\cite{demand2023} & LNCS (MIWAI), 2023 & Price, demand \\ \hline
S17 & Point and interval forecasting, decomposition-ensemble and KDE~\cite{kde2024} & Soft Comput., 2024 & Price, interval \\ \hline
S18 & Optimised machine learning for commodity prices~\cite{optimized2023} & Neural Comput. Appl., 2024 & Price \\ \hline
S19 & GCN--BiGRU--TPE futures price prediction~\cite{gcn2024} & Comput. Econ., 2025 & Price, graph \\ \hline
S20 & Crop yield prediction and recommendation, deep ensembles (survey)~\cite{cropyield} & IEEE conf., 2024 & Crop, survey \\ \hline
S21 & Framework for ragi price prediction in Karnataka~\cite{meena2024} & IEEE Access, 2024 & Price \\ \hline
S22 & AgTech: volatility prediction for commodity trading~\cite{agtech2024} & IEEE Access, 2024 & Volatility \\ \hline
S23 & Smart agriculture: state, opportunities and challenges~\cite{smartag2024} & IEEE Access, 2024 & Review \\ \hline
S24 & Smart farming: crop recommendation and price prediction~\cite{smartfarm2025} & IEEE conf., 2025 & Crop, price \\ \hline
S25 & Federated explainable AI for smart agriculture~\cite{fedxai2025} & IEEE Access, 2025 & Crop, XAI \\ \hline
\end{longtable}}

\subsection{Techniques and results}
Tables~\ref{tab:survey-t1} to~\ref{tab:survey-t3} record, for each study, the techniques used and the principal results as reported by the authors.

\Needspace{12\baselineskip}
{\footnotesize\setlength{\tabcolsep}{4pt}\renewcommand{\arraystretch}{1.15}
\begin{longtable}{|p{0.9cm}|p{5.8cm}|p{9.2cm}|}
\caption{Techniques and results, S1--S10 (price, volatility and anomaly studies, Elsevier)}\label{tab:survey-t1}\\
\hline \textbf{ID} & \textbf{Techniques} & \textbf{Principal results} \\ \hline \endfirsthead
\hline \textbf{ID} & \textbf{Techniques} & \textbf{Principal results} \\ \hline \endhead
\hline \endfoot
S1 & ARIMA, VAR, LSTM and stochastic simulation (geometric, fractional and bridge Brownian motion); value-at-risk & LSTM best in high-volatility and crisis periods; ARIMA stable in short-term, low-volatility markets; geometric Brownian motion balances accuracy and risk; fractional motion captures long memory but raises volatility \\ \hline
S2 & Transformer forecasting; attention-boosted LSTM variational autoencoder for anomalies & Transformers beat LSTM and GRU in price forecasting; the attention-boosted autoencoder improves anomaly detection and captures sudden spikes and drops \\ \hline
S3 & ARIMA plus LSTM hybrid; random-forest selection of lags and features; volatility modelling & ARIMA captures linear trend, LSTM the non-linear part, random forest improves feature selection; the hybrid cuts RMSE and MAPE markedly \\ \hline
S4 & ICEEMDAN and VMD decomposition followed by GRU & Decomposition reduces noise and non-stationarity; the hybrid outperforms LSTM and RNN variants \\ \hline
S5 & Heterogeneous autoregressive (HAR) models, LASSO; news-based and market-based uncertainty indicators & Market-based indicators outperform news-based ones; LASSO improves on HAR; uncertainty is strongly linked to price volatility \\ \hline
S6 & HAR-RV variants on high-frequency realised volatility & Financial stress predicts realised volatility in-sample and shows financialisation of agricultural markets; out-of-sample prediction is inconsistent \\ \hline
S7 & EWT, SSA and VMD decomposition; ARIMA, ETS, BPNN, ELM; PSO-CS weight optimisation & The optimised ensemble beats all single models, with a large reduction in MAPE \\ \hline
S8 & Variable temporal transformer with self-attention for multivariate anomaly detection & Captures temporal and inter-variable dependence, beats traditional detectors and gives interpretable anomaly explanations \\ \hline
S9 & Hidden Markov model with RNN, LSTM, GRU, BiLSTM, CNN; technical indicators and FFT & The Markov model captures regime change; guided deep models beat the same models alone for short and medium horizons \\ \hline
S10 & Time-varying-parameter VAR for transmission of policy uncertainty & Policy uncertainty affects price volatility, with effects that vary over time and across commodities \\ \hline
\end{longtable}}

\Needspace{12\baselineskip}
{\footnotesize\setlength{\tabcolsep}{4pt}\renewcommand{\arraystretch}{1.15}
\begin{longtable}{|p{0.9cm}|p{5.8cm}|p{9.2cm}|}
\caption{Techniques and results, S11--S19 (Springer Nature)}\label{tab:survey-t2}\\
\hline \textbf{ID} & \textbf{Techniques} & \textbf{Principal results} \\ \hline \endfirsthead
\hline \textbf{ID} & \textbf{Techniques} & \textbf{Principal results} \\ \hline \endhead
\hline \endfoot
S11 & STL decomposition into trend, seasonal and remainder; an LSTM per component; recombined forecast & Beat ARIMA and a plain LSTM; seasonal components predicted more reliably; RMSE and MAPE improved on benchmarks \\ \hline
S12 & CEEMDAN decomposition with a time-delay neural network per component; compared with ARIMA, random forest, GBM, SVR & Best of all models; RMSE about 20\% lower than other EMD-based models and about 57\% lower than traditional machine learning; MAPE near 2--4\% for major crops \\ \hline
S13 & NBEATSX and TransformerX with weather variables; compared with ARIMAX, ANN, SVR, random forest, XGBoost; RMSE, MAE, sMAPE, MASE & NBEATSX best for all three crops (tomato, onion, potato); deep models clearly beat statistical and machine-learning baselines \\ \hline
S14 & LSTM for long-term pattern; GARCH on LSTM residuals; compared with GARCH, GJR-GARCH, EWMA, MEM & Best for all commodities; onion RMSE fell from 5.44 (GARCH) to 0.70, tomato from 7.67 to 0.79 \\ \hline
S15 & ARIMA, AR, MA, ES, STL for yield; supply from yield, residue and imports; demand by time series; decision tree, RF, KNN, kernel ridge for price & ARIMA(4,1,1) best for yield; decision-tree regressor best for price (RMSE about 0.115, MAPE about 14\%) \\ \hline
S16 & RNN, LSTM and bidirectional LSTM compared with statistical and machine-learning methods; horizon up to one year & LSTM and Bi-LSTM capture temporal dependence better than ARIMA variants; detailed error figures are not available in the accessible preview \\ \hline
S17 & VMD with energy entropy; K-means reconstruction; ARIMA for low and LSTM for high frequencies; WOA-optimised adaptive KDE intervals & Lowest error among compared methods; MAPE about 1--2\% one step ahead; interval coverage (PICP) above 95\% at the 95\% level; 5--15\% lower RMSE than other hybrids \\ \hline
S18 & Extreme learning machine and LSTM tuned by genetic algorithm; Auto-ARIMA; 11 commodities, daily 2000--2022 & GA-ELM best overall (MAPE about 1--3\%), trains faster than GA-LSTM, forecasts 60 days ahead and follows the 2022 peaks \\ \hline
S19 & Multi-graph construction; graph convolution for spatial structure; BiGRU for time; TPE Bayesian tuning & On corn futures, RMSE about 20.2, MAPE about 0.55\%, $R^2$ about 0.97 \\ \hline
\end{longtable}}

\Needspace{12\baselineskip}
{\footnotesize\setlength{\tabcolsep}{4pt}\renewcommand{\arraystretch}{1.15}
\begin{longtable}{|p{0.9cm}|p{5.8cm}|p{9.2cm}|}
\caption{Techniques and results, S20--S25 (IEEE)}\label{tab:survey-t3}\\
\hline \textbf{ID} & \textbf{Techniques} & \textbf{Principal results} \\ \hline \endfirsthead
\hline \textbf{ID} & \textbf{Techniques} & \textbf{Principal results} \\ \hline \endhead
\hline \endfoot
S20 & PSO and WOA for feature selection; SVM, random forest, RNN, Bi-LSTM, CNN, MLP and ensembles (survey of methods) & Ensemble and optimisation approaches outperform traditional methods in crop prediction; Bi-LSTM and CNN strongest on climate and crop data \\ \hline
S21 & LSTM deep-learning framework with temporal and seasonal features; compared with statistical and machine-learning models & Higher accuracy on ragi price trends, robust to non-linear fluctuations, better precision and recall than baselines \\ \hline
S22 & Bi-LSTM on multivariate market indicators (price, volume, open interest); rolling-window volatility index; dropout; binary classification of high/low volatility & Improved accuracy of volatility forecasts for exchange-traded commodities; robust to sudden fluctuations; fewer false predictions \\ \hline
S23 & Survey of IoT, AI and machine learning, blockchain, cloud and precision farming & Documents gains in resource efficiency (water, fertiliser, energy) and applications in precision farming, monitoring and yield prediction \\ \hline
S24 & XGBoost, random forest, decision tree, SVM, KNN, naive Bayes, linear and logistic regression & Integrated crop recommendation and price prediction; ensembles most accurate; price models more reliable than traditional statistics \\ \hline
S25 & Federated learning with explainable AI; CNN, LSTM, RF, SVM, gradient boosting; SHAP and LIME & Better crop-recommendation accuracy across distributed data while preserving privacy; transparent predictions; scales across regions \\ \hline
\end{longtable}}

\subsection{Research gaps and takeaways}
Tables~\ref{tab:survey-g1} to~\ref{tab:survey-g3} record the gaps named for each study in the workbook and the lesson each gives for MandiSense AI.

\Needspace{12\baselineskip}
{\footnotesize\setlength{\tabcolsep}{4pt}\renewcommand{\arraystretch}{1.15}
\begin{longtable}{|p{0.9cm}|p{6.4cm}|p{8.6cm}|}
\caption{Research gaps and takeaways, S1--S10}\label{tab:survey-g1}\\
\hline \textbf{ID} & \textbf{Research gaps} & \textbf{Takeaways for this project} \\ \hline \endfirsthead
\hline \textbf{ID} & \textbf{Research gaps} & \textbf{Takeaways for this project} \\ \hline \endhead
\hline \endfoot
S1 & Models used independently; no shock detection; volatility not in decision support; no explainable architecture & No single model wins in every regime, so a council of specialists is justified; risk-aware, adaptive systems are needed \\ \hline
S2 & Centralised, not agent-based; no explicit risk quantification; limited decision-support integration & Joint forecasting and anomaly detection works; risk and explanation still need to be added \\ \hline
S3 & Monthly data only; no ensemble extension; volatility handled indirectly & Hybrid statistical-plus-deep models beat single models; feature selection matters \\ \hline
S4 & High computational cost; no volatility or anomaly module; one pipeline & Noise reduction is critical for agricultural series; decomposition plus deep learning is effective (and, here, audited for leakage) \\ \hline
S5 & Volatility only, not price; not linked to forecasting or decisions & External uncertainty signals matter; volatility should be tied to price prediction \\ \hline
S6 & Out-of-sample reliability weak; not linked to price models & Financial stress is a volatility driver; regimes change; adaptive models are needed \\ \hline
S7 & Centralised; no explainability or volatility alerts; costly & Combination improves accuracy and stability; weight optimisation is crucial \\ \hline
S8 & Not integrated with price forecasting; little agricultural use; intensive & Attention helps anomaly detection; anomaly detection should be coupled with forecasting \\ \hline
S9 & Commodity- and region-specific; no ensemble or multi-agent design; limited explainability & Regime awareness improves forecasts; probabilistic-plus-deep hybrids suit volatile series \\ \hline
S10 & China only; no forecasting framework & Policy uncertainty drives volatility; time-varying models are needed; policy effects belong in forecasting systems \\ \hline
\end{longtable}}

\Needspace{12\baselineskip}
{\footnotesize\setlength{\tabcolsep}{4pt}\renewcommand{\arraystretch}{1.15}
\begin{longtable}{|p{0.9cm}|p{6.4cm}|p{8.6cm}|}
\caption{Research gaps and takeaways, S11--S19}\label{tab:survey-g2}\\
\hline \textbf{ID} & \textbf{Research gaps} & \textbf{Takeaways for this project} \\ \hline \endfirsthead
\hline \textbf{ID} & \textbf{Research gaps} & \textbf{Takeaways for this project} \\ \hline \endhead
\hline \endfoot
S11 & Weather, supply and demand shocks and market events not used; centralised & Decomposition handles strong seasonality, as in mandi series; preprocessing to isolate components helps deep models \\ \hline
S12 & History only, no weather or demand; centralised & Decomposition before prediction helps volatile prices; hybrids are more reliable than single models \\ \hline
S13 & Only weather; no demand, festivals or exports; centralised & External factors improve accuracy; multi-input forecasting suits Indian markets \\ \hline
S14 & Price history only; centralised & LSTM--GARCH suits abnormal-volatility detection and risk alerts; volatility should join multi-model forecasting \\ \hline
S15 & One crop (corn); no deep learning or real-time data & Integrating supply and demand improves price forecasts; multi-factor beats price alone \\ \hline
S16 & Few metrics reported; no exogenous drivers & Deep sequence models capture temporal trend; test several architectures; report detailed errors \\ \hline
S17 & Price history only; computationally complex & Intervals support risk-aware decisions; ARIMA plus LSTM improves stability; ranges should accompany forecasts \\ \hline
S18 & No external factors; centralised & Optimised hybrids are stable; 60-day forecasts are feasible; compare several models; add arrivals, season and festivals \\ \hline
S19 & One commodity; no weather or demand cycles & Graph structure suits mandis linked by trade; Bayesian tuning automates configuration \\ \hline
\end{longtable}}

\Needspace{12\baselineskip}
{\footnotesize\setlength{\tabcolsep}{4pt}\renewcommand{\arraystretch}{1.15}
\begin{longtable}{|p{0.9cm}|p{6.4cm}|p{8.6cm}|}
\caption{Research gaps and takeaways, S20--S25}\label{tab:survey-g3}\\
\hline \textbf{ID} & \textbf{Research gaps} & \textbf{Takeaways for this project} \\ \hline \endfirsthead
\hline \textbf{ID} & \textbf{Research gaps} & \textbf{Takeaways for this project} \\ \hline \endhead
\hline \endfoot
S20 & Few multi-source datasets; no standard feature engineering; narrow metrics; little interpretability & Ensembles and optimisation help; precision, recall and F1 should be reported; false rates affect trust \\ \hline
S21 & One crop; no socio-economic factors; no deployment study; interpretability & LSTM models suit price series; temporal features help; the framework supports farmers and policymakers \\ \hline
S22 & Policy factors missing; no deployment or multi-commodity scale; interpretability; cross-market links & Volatility prediction can steady trading decisions; ensembles beat single models; external features help \\ \hline
S23 & Little real-time deployment in developing regions; interoperability; privacy; farmer usability & Smart agriculture reduces waste and raises productivity; adoption depends on farmers; solutions must suit smallholders \\ \hline
S24 & No real-time integrated deployment; socio-economic and policy factors; interpretability; scalability & Crop recommendation and price forecasting together support market decisions; multi-source data is crucial \\ \hline
S25 & No real-time rural deployment; interoperability; efficiency against explainability; farmer usability; security & Federated learning avoids sharing raw farmer data; explanations build trust \\ \hline
\end{longtable}}

\subsection{Comparison matrix}
Table~\ref{tab:survey-matrix} compares the studies on ten properties. A tick means the study uses or addresses the property. The last row shows MandiSense AI for reference.

\begin{table}[!t]
\centering
\caption{Comparison of the surveyed studies on ten properties}
\label{tab:survey-matrix}
\footnotesize
\setlength{\tabcolsep}{4.6pt}
\renewcommand{\arraystretch}{1.1}
\begin{tabular}{|l|c|c|c|c|c|c|c|c|c|c|}
\hline
\textbf{ID} & \textbf{Decomp.} & \textbf{RNN} & \textbf{Attn.} & \textbf{Volat.} & \textbf{Ensemble} & \textbf{Exog.} & \textbf{Interval} & \textbf{Anomaly} & \textbf{Graph} & \textbf{Decision} \\ \hline
S1  & & $\checkmark$ & & $\checkmark$ & $\checkmark$ & & & & & \\ \hline
S2  & & & $\checkmark$ & & & & & $\checkmark$ & & \\ \hline
S3  & & $\checkmark$ & & $\checkmark$ & $\checkmark$ & & & & & \\ \hline
S4  & $\checkmark$ & $\checkmark$ & & & & & & & & \\ \hline
S5  & & & & $\checkmark$ & & $\checkmark$ & & & & \\ \hline
S6  & & & & $\checkmark$ & & $\checkmark$ & & & & \\ \hline
S7  & $\checkmark$ & & & & $\checkmark$ & & & & & \\ \hline
S8  & & & $\checkmark$ & & & & & $\checkmark$ & & \\ \hline
S9  & & $\checkmark$ & & $\checkmark$ & & & & & & \\ \hline
S10 & & & & $\checkmark$ & & $\checkmark$ & & & & \\ \hline
S11 & $\checkmark$ & $\checkmark$ & & & & & & & & \\ \hline
S12 & $\checkmark$ & & & & & & & & & \\ \hline
S13 & & & $\checkmark$ & & & $\checkmark$ & & & & \\ \hline
S14 & & $\checkmark$ & & $\checkmark$ & & & & & & \\ \hline
S15 & & & & & & $\checkmark$ & & & & \\ \hline
S16 & & $\checkmark$ & & & & & & & & \\ \hline
S17 & $\checkmark$ & $\checkmark$ & & & $\checkmark$ & & $\checkmark$ & & & \\ \hline
S18 & & $\checkmark$ & & & $\checkmark$ & & & & & \\ \hline
S19 & & $\checkmark$ & & & & & & & $\checkmark$ & \\ \hline
S20 & & $\checkmark$ & & & $\checkmark$ & $\checkmark$ & & & & $\checkmark$ \\ \hline
S21 & & $\checkmark$ & & & & & & & & \\ \hline
S22 & & $\checkmark$ & & $\checkmark$ & & $\checkmark$ & & & & \\ \hline
S23 & & & & & & & & & & \\ \hline
S24 & & & & & $\checkmark$ & $\checkmark$ & & & & $\checkmark$ \\ \hline
S25 & & $\checkmark$ & & & & & & & & $\checkmark$ \\ \hline
\textbf{Total} & \textbf{5} & \textbf{14} & \textbf{3} & \textbf{8} & \textbf{7} & \textbf{8} & \textbf{1} & \textbf{2} & \textbf{1} & \textbf{3} \\ \hline
\textbf{This work} & $\checkmark$ & & $\checkmark$ & $\checkmark$ & $\checkmark$ & & $\checkmark$ & $\checkmark$ & $\checkmark$ & $\checkmark$ \\ \hline
\end{tabular}
\par\smallskip{\footnotesize Decomp.: decomposition; Attn.: attention or transformer; Volat.: volatility or risk modelling; Exog.: external drivers; Interval: prediction intervals; Graph: spatial or graph structure; Decision: output tied to a farmer decision. \emph{This work}: leakage audit of decomposition, zero-shot transformer baseline, volatility desk, council, conformal floors, shock desk, spillover engine, market choice.}
\end{table}

\subsection{Gap analysis across the survey}
The gaps recorded for the 25 studies fall into a small number of themes. Table~\ref{tab:survey-gaps} counts how many studies name each and states how the present work responds.

\begin{table}[!t]
\centering
\caption{Recurring gaps in the surveyed studies and the response of this project}
\label{tab:survey-gaps}
\small
\renewcommand{\arraystretch}{1.15}
\begin{tabular}{|p{4.0cm}|p{3.4cm}|c|p{6.4cm}|}
\hline
\textbf{Gap theme} & \textbf{Studies} & \textbf{n} & \textbf{Response in MandiSense AI} \\ \hline
External drivers (weather, demand, policy, socio-economic) not used & S11--S19, S21, S22, S24 & 12 & Arrivals, season and calendar features in the panel; a decision-focused label; no weather feed (live data not available) \\ \hline
Centralised design, no multi-agent or distributed structure & S2, S3, S7, S9, S11, S12, S14, S18 & 8 & Specialist council with learned weights and per-crop engines behind one API (Section~\ref{sec:architecture}) \\ \hline
Weak interpretability or explainability & S1, S7, S9, S20--S22, S24, S25 & 8 & Why-this-mandi breakdown, what-if, track record and a checked brief (Section~\ref{sec:insight-method}) \\ \hline
No real-time deployment or scalability study & S5, S6, S10, S21--S25 & 8 & A running service, a 60-check smoke test and a morning digest; no field trial yet \\ \hline
Volatility, risk or anomaly not tied to price or decisions & S1, S2, S4, S5, S6, S8, S10 & 7 & Shock desk, volatility view and revenue floors in the same service as the advice \\ \hline
Narrow data (one crop, market or region; monthly only) & S3, S9, S10, S15, S19, S21 & 6 & Four crops, 12 markets, daily data; one region (stated as a limit) \\ \hline
High computational cost & S4, S7, S8, S17 & 4 & Gradient boosting served; heavy models used only as challengers \\ \hline
Judged by forecast error, not by a decision & S1--S4, S7, S9, S11--S19, S21 & 16 & Regret and floor violation as the primary measures (Chapter~\ref{ch:results}) \\ \hline
Decomposition leakage not audited (our reading) & S4, S7, S11, S12, S17 & 5 & Walk-forward re-decomposition audit (Section~\ref{sec:leakage}) \\ \hline
\end{tabular}
\end{table}

Three conclusions follow. First, hybrid and decomposition-based models dominate recent work (14 of 25 use recurrent networks, 5 use decomposition) and are judged almost entirely by forecast error. Second, only one study (S17) provides an interval and none states a guarantee or a downside in rupees. Third, no study treats the farmer's own choice, of market and day, as the quantity to be optimised, and no study audits how much of the reported gain comes from fitting a decomposition on the whole series. These are the gaps that Section~\ref{sec:limitations} turns into objectives.


\section{Decisions, Not Forecasts}
\label{sec:ml-decision}

When forecasts feed an optimisation, minimising forecast error is not the same as minimising decision loss. The smart ``predict, then optimise'' (SPO+) loss trains a model directly on the regret of the decision it induces~\cite{elmachtoub2022}, and a recent benchmark across many decision-focused methods found the benefit to depend heavily on the problem~\cite{mandi2024}. For market choice the question is open, and Section~\ref{sec:effects} answers it for this problem. The neural variant used here encodes the set of markets with a permutation-invariant network~\cite{zaheer2017}, and a robust loss limits the influence of outlying prices on training~\cite{huber1964}.

\section{Calibrated Uncertainty}
\label{sec:uncertainty}

Conformal prediction gives intervals with finite-sample coverage under exchangeability~\cite{vovk2005,angelopoulos2023}, and conformalised quantile regression adapts the width to the input~\cite{romano2019}. Prices drift, so exchangeability fails; adaptive conformal inference tracks a long-run miscoverage rate under shift~\cite{gibbs2021}. A further complication, that the outcome of a sale is observed only days later, is handled in Section~\ref{sec:floors}.

\section{Missing and Informatively Missing Data}
\label{sec:missing}

A price is observed only when a market reports. If reporting depends on conditions the model also cares about, the data are not missing at random~\cite{rubin1976}, and learning from observed values alone is biased. Neural models for multivariate series with missing values add the missingness pattern itself as an input~\cite{che2018}; the same idea appears here as a reporting head and as silence features (Section~\ref{sec:formulation}).

\section{Transmission Between Commodities and the Effect of Followers}
\label{sec:transmission}

Prices of crops that share land, labour, storage or consumers move together, and some of that movement is the common market factor and the common season rather than transmission from one crop to another. Studies of transmission typically estimate it from whole series, which blends the three. The method used here removes the first common factor and each crop's calendar-month mean, defines shock episodes, and compares the response after an episode with the unconditional baseline, resampling over episodes rather than days (Section~\ref{sec:spill-method}) because days within an episode are strongly dependent~\cite{kunsch1989}; many pairs are tested, so multiplicity is controlled~\cite{holm1979}. Where the evidence is thin, the system withholds the claim.

A recommendation served to many users also changes what it predicts: a mandi that receives more produce than usual pays less, so advice that is good for one farmer can be poor for a thousand. This thesis treats that effect as a simulation with a stated price flexibility, not as a measured law (Section~\ref{sec:insight-method}). We did not find prior work that evaluates it for mandi choice, and cite none.

\section{Language Models for Explanation}
\label{sec:llm}

Language models can turn model output into readable advice in several languages, and tool-calling interfaces let a model fill a fixed schema instead of writing free text. The risk is that the text is fluent whether or not it is faithful: a model asked about prices will produce prices, some plausible and wrong. Work on foundation models for agricultural forecasting concerns the forecast itself~\cite{ansari2024,wang2026}; among the studies surveyed we found none that couples a language model to a price decision with a check that every stated figure appears in the evidence. This project treats the language model as a writer whose output is verified and may be discarded (Section~\ref{sec:intel-method}). The writer is an open-weight mixture-of-experts model~\cite{gptoss2025} and the judge is a different model family~\cite{qwen3}, so that the writer does not grade itself. Anthropic's Claude~\cite{claude3} was used by the authors as a development and writing assistant; Section~\ref{sec:ai-disclosure} describes that use.

\section{Limitations of Existing Systems}
\label{sec:limitations}

Table~\ref{tab:gaps} lists the limitations found and how each is addressed here. The first eight are properties of the literature; the last two are properties of what farmers actually have today: a price portal without advice, and no message that reaches them.

\begin{table}[!t]
    \centering
    \caption{Gaps in existing approaches and the objective that addresses each}
    \label{tab:gaps}
    \small
    \begin{tabular}{|p{4.6cm}|p{6.6cm}|p{3cm}|}
    \hline
    \textbf{Limitation} & \textbf{Description and impact} & \textbf{Addressed by} \\ \hline
    Judged on forecast error & Models are ranked by RMSE or MAPE; a model can win on error and still pick the wrong market~\cite{elmachtoub2022,mandi2024} & Objectives 2, 3 (Sec.~\ref{sec:effects}) \\ \hline
    Reporting gaps ignored & Training on posted prices only teaches the model that every market is open~\cite{rubin1976,che2018} & Objective 2 (Sec.~\ref{sec:formulation}) \\ \hline
    No statement of downside & Point forecasts, and intervals without guarantees, give a farmer nothing firm to set against the cost of a bad sale~\cite{kde2024,gibbs2021} & Objective 4 (Sec.~\ref{sec:floors}) \\ \hline
    Leakage in hybrid evaluation & Decompositions fitted on the full series flatter the reported skill~\cite{jaiswal2025,sathe2025,bergmeir2012} & Objective 6 (Sec.~\ref{sec:leakage}) \\ \hline
    Language-model advice is unchecked & A model can state a price that is nowhere in the evidence, and nothing in the text shows it & Objective 7 (Sec.~\ref{sec:intel-method}) \\ \hline
    Cross-crop effects confounded & Common movement and season look like transmission when estimated on whole series & Objective 8 (Sec.~\ref{sec:spill-method}) \\ \hline
    Advice cannot be inspected & A recommendation without its cost arithmetic, the cost at which it changes and its record cannot be checked by a farmer & Objective 9 (Sec.~\ref{sec:insight-method}) \\ \hline
    Followers ignored & Advice built for one farmer is silent on what happens when many take it & Objective 9 (Sec.~\ref{sec:insight-method}) \\ \hline
    Portal shows prices, not advice & Yesterday's quote at each market; no transport, shrinkage or timing~\cite{agmarknet} & Objective 5 (Ch.~\ref{ch:impl}) \\ \hline
    No message reaches the farmer & Information sits on a portal; nothing arrives in the morning, in the farmer's language, with the date of every price~\cite{fafchamps2012} & Objective 10 (Sec.~\ref{sec:digest-method}) \\ \hline
    \end{tabular}
\end{table}

\section{Problem Statement}
\label{sec:problem}

The literature supports the following reading. Agricultural price forecasting is mature and mostly judged by forecast error~\cite{wu2026}; decision-focused learning~\cite{elmachtoub2022,mandi2024} and calibrated uncertainty~\cite{gibbs2021,angelopoulos2023} are mature elsewhere; and information helps farmers when it changes a specific decision~\cite{jensen2007,aker2010,fafchamps2012}. These strands have not been joined on real mandi data with the missing-report problem, an audit of reported gains, a check on language-model advice, and a channel that reaches the farmer.

\paragraph{Problem.} Given the daily, partly missing prices reported by the mandis around Bengaluru and Kolar, the cost of moving a crop and its perishability, recommend to a farmer \emph{where and when to sell}, state a revenue floor that holds at a stated rate on later days, explain the recommendation in checkable terms in the farmer's language, and report honestly where the method does not beat simple alternatives. Chapter~\ref{ch:method} defines this precisely, Chapter~\ref{ch:impl} builds it and Chapter~\ref{ch:results} measures it.

\paragraph{Scope of the claim.} The evidence is offline, on archived data from one region, with no live deployment and no field trial. Where the literature reports larger effects, for example the welfare gains of information interventions~\cite{jensen2007,aker2010}, our results do not test them.
